% Глава: 01-intro
% Статус: наполнено (первый прогон).
\chapter{Введение}
\label{ch:01-intro}

\section{О дисциплине}

Курс \emph{``Введение в операционные системы и инструменты разработчика''}
закладывает базовые навыки, которые потом используются на курсах
\emph{Объектно-ориентированное программирование},
\emph{Программная архитектура} и \emph{Операционные системы}.
В отличие от курса по программированию, здесь главный предмет --- не язык,
а \textbf{среда, в которой вы пишете и запускаете программы}:
командная оболочка, файловая система, процессы, компиляция, сборка,
взаимодействие процессов и виртуализация.

Цель дисциплины --- познакомить с основными инструментами разработки:
командными интерпретаторами и основами скриптинга, удалённым доступом,
системами управления версиями, сборщиками и отладчиками. По итогам курса
вы должны \emph{уметь} читать системную документацию, разбираться в том,
что делает система под капотом, и писать небольшие программы на C с их
сборкой, тестированием и отладкой.

Курс читается на первом курсе (оба семестра), без экзамена.
Контроль состоит из трёх частей:
\begin{itemize}
  \item \textbf{портфолио} --- накопленный набор выполненного: работа в
    терминале, собственные сценарии и программы, история коммитов;
  \item \textbf{групповой/итоговый проект} --- небольшое совместное изделие,
    собранный в единый репозиторий;
  \item \textbf{дифференцированный зачёт} --- один в каждом семестре,
    оценка по шкале \emph{неуд / уд / хорошо / отлично}.
\end{itemize}

\section{Рабочая среда курса: сервер лаборатории}

Основное рабочее место на курсе --- не ваш компьютер, а \textbf{сервер
лаборатории} \fpath{pglab.nsu.ru}. Курсу выделено отдельное
вычислительное устройство: на нём разворачивается по \textbf{одной
персональной виртуальной машине (ВМ) на каждого студента} --- изолированная
система на базе Ubuntu, доступная только по \textbf{SSH}. Развёртывание и
поддержание такого окружения --- самостоятельная инфраструктурная задача
университета; характеристики ВМ и настройка доступа --- в
\nameref{ch:02-linux-access}.

Вся практическая работа дисциплины выполняется именно на этом сервере:
работа в терминале и с файловой системой; написание, сборка, тестирование
и отладка программ на C (главы~\ref{ch:13-c-compilation}--\ref{ch:17-debugging});
контейнеры Docker (главы~\ref{ch:19-docker-basics}--\ref{ch:23-docker-resources});
Git и итоговый проект (\nameref{ch:25-git}).

Доступ к своей ВМ открыт в любое время: из кампуса --- напрямую, извне ---
через университетский \textbf{VPN} (см.\ \nameref{ch:02-linux-access}). По сути,
сервер лаборатории --- это «виртуальный компьютер» каждого студента на весь
семестр. Нагрузка на него существенна и устойчива: одновременно работающая
группа, у каждого --- своя ВМ, свои процессы, компиляции и контейнеры.

\begin{tipbox}
Локальный Linux на личном компьютере не требуется: вся среда --- на
\fpath{pglab.nsu.ru}. Подключаемся по SSH: \cmd{ssh <login>@pglab.nsu.ru}
(данные для входа выдаёт семинарист).
\end{tipbox}

\section{Как устроена методичка}

Методичка \textbf{тематическая}, а не ``по номерам семинаров''.
Одна тема может покрываться несколькими семинарами, а один семинар
--- несколькими темами. Главы собраны в пять частей:

\begin{center}
\begin{tabular}{ll}
\toprule
\textbf{Часть} & \textbf{Содержание} \\
\midrule
Часть 0  & Введение и рабочая среда \\
Часть I  & Операционная система и командная строка \\
Часть II & Разработка на C и сборка \\
Часть III& Взаимодействие процессов и виртуализация (Docker) \\
Часть IV & Среда разработчика (Git, CI, скриптинг) \\
\bottomrule
\end{tabular}
\end{center}

Каждая глава --- самостоятельный блок. Внутри главы есть
\emph{связь с работой ОС} (что происходит под капотом),
\emph{типичные ошибки}, \emph{ссылки на документацию}
(главным образом \cmd{man}) и \emph{задачи для самопроверки}.
Перекрёстные ссылки вида «см.\ главу~\ref{ch:03-shell}» ведут к связанным
темам. Состояние наполнения глав и их связь с лекциями и лабораторными
показывает \nameref{app:roadmap}.

\section{Как пользоваться методичкой}

\begin{itemize}
  \item Команды, флаги, переменные и термины написаны на английском
    (\cmd{ls -la}, \fpath{HOME}, \emph{container}) --- так их проще найти в
    документации и в интернете.
  \item Путь к файлу или директории оформлен как \fpath{/etc/passwd}.
  \item \emph{Задачи для самопроверки} в конце каждой главы --- это
    минимальный набор действий, которым вы должны владеть. Если что-то из них
    не удаётся выполнить --- перечитайте главу.
  \item Запоминать наизусть все флаги \textbf{не нужно}. Навык, который мы
    отрабатываем, --- быстро находить нужный параметр в документации
    (\nameref{ch:06-man}).
\end{itemize}

\begin{warnbox}
Если задача на семинаре не успевает быть решена, она автоматически
становится домашней. Методичка даёт материал и самопроверку, но выполнение
самостоятельно --- ваша работа.
\end{warnbox}

\section{Материалы и дальнейшее развитие}

Методичка растёт по мере прохождения курса: новые темы добавляются
по мере того, как появляются соответствующие семинары и лекции.
Главы, по которым материал ещё не развёрнут, помечены как
\emph{будет дополнено} и видны в \nameref{app:roadmap}.
Шпаргалки по командам собраны в \nameref{app:cheatsheets},
определения терминов --- в \nameref{app:glossary},
основная литература и ссылки --- в \nameref{app:references}.

\section{Задачи для самопроверки}
\begin{enumerate}
  \item Назовите, на какие последующие дисциплины опирается данный курс
    и какая у него основная цель.
  \item Какие три формы контроля составляют итоговую оценку и чем
    дифференцированный зачёт отличается от экзамена?
  \item Назовите пять частей методички и скажите, какой типичной теме
    посвящена каждая.
  \item Объясните, почему на курсе не требуется заучивать флаги команд
    наизусть и какой навык важнее.
\end{enumerate}
